% ============================================================
% MISE EN PAGE — VARIANTE « SIDEBAR » (deux colonnes)
% ============================================================
% Marges quasi nulles : le fond coloré de la sidebar est un aplat qui part du
% bord gauche de la page, le contenu doit donc pouvoir s'en approcher.
\geometry{left=0.4cm, right=0.4cm, top=0cm, bottom=0.1cm}

% ------------------------------------------------------------
% DENSITÉ DU CONTENU
% ------------------------------------------------------------
% La colonne principale ne fait que 0,75\textwidth : les sous-puces de détail
% y débordaient sur deux lignes chacune, ce qui poussait la fin du CV sur une
% seconde page. Or les deux minipages (sidebar + contenu) forment UNE ligne
% horizontale insécable : dès qu'elle dépasse la hauteur de page, LaTeX la
% reporte en bloc et laisse la page 1 entièrement vide. Il ne s'agit donc pas
% d'un débordement cosmétique — c'est ce qui cassait le PDF.
\renewcommand{\cvsub}[1]{{\footnotesize\color{detailcolor}\itshape #1}}

% Interlignage resserré, même parti pris que la variante une colonne.
\linespread{0.98}\selectfont
